% chapter2_theorie.tex
% Tweede hoofdstuk - Theoretische achtergrond
% Demonstreert formules, symbolen en siunitx

\chapter{Theoretische Achtergrond}
\label{ch:theorie}

Dit hoofdstuk demonstreert het gebruik van formules, symbolen en eenheden.

\section{Symbolen en Eenheden}

Bij de eerste vermelding van een symbool gebruiken we \verb|\sym{}|:

\sym{v}{Snelheid}{m/s}

Later in de tekst verschijnt alleen het symbool: de \sym{v}{Snelheid}{m/s} van het object.

\subsection{Eenheden met siunitx}

\begin{theorieblok}[Correcte eenheden met siunitx]
Gebruik \textbf{altijd} \verb|\SI{}{}| voor getallen met eenheden:
\begin{itemize}
    \item Snelheid: \verb|\SI{50}{km/h}| → \SI{50}{km/h}
    \item Temperatuur: \verb|\SI{20}{\degreeCelsius}| → \SI{20}{\degreeCelsius}
    \item Kracht: \verb|\SI{1.5}{kN}| → \SI{1.5}{kN}
    \item Druk: \verb|\SI{101.325}{kPa}| → \SI{101.325}{kPa}
\end{itemize}

Voor getallen zonder eenheid: \verb|\num{1234567}| → \num{1234567}

Voor hoeken: \verb|\ang{45}| → \ang{45}
\end{theorieblok}

\section{Formules}

Elke belangrijke formule wordt geregistreerd met \verb|\frm{}{}{}|:

\frm{Kinetische energie}{E_k = \frac{1}{2}mv^2}{waarbij $m$ de massa [\si{kg}] en $v$ de snelheid [\si{m/s}]}

\frm{Potentiële energie}{E_p = mgh}{met $g = \SI{9.81}{m/s^2}$ en $h$ de hoogte [\si{m}]}

\section{Wiskunde shortcuts}

\subsection{Afgeleiden}

De afgeleide notatie is vereenvoudigd.

Voor een enkele formule gebruik je \verb|\[ ... \]|:
\[ \dv{f}{x} = \lim_{\Delta x \to 0} \frac{\Delta f}{\Delta x} \]

Voor meerdere formules die je mooi onder elkaar wilt uitlijnen (op de $=$) kun je \verb|\[ ... \]| gebruiken met meerdere displays:
\[ \dd{y}{x} = \text{gewone afgeleide} \]
\[ \pd{f}{x} = \text{partiële afgeleide} \]
\[ \ddn{y}{x}{2} = \text{tweede afgeleide} \]

of met \verb|\begin{aligned} ... \end{aligned}| je creert een soort boxomgeving voor wiskundige uitlijning:
\[
\begin{aligned}
  \dd{y}{x} &= \text{gewone afgeleide} \\
  \pd{f}{x}  &= \text{partiële afgeleide}
\end{aligned}
\]


\subsection{Vectoren en verzamelingen}

\begin{itemize}
    \item Vector: $\vec{F} = m\vec{a}$
    \item Eenheidsvector: $\uvec{n}$
    \item Reële getallen: $x \in \RR$
    \item Complexe getallen: $z \in \CC$
\end{itemize}

\section{Cross-references met cleveref}

Gebruik \verb|\cref{}| voor automatische Nederlandse referenties:
\begin{itemize}
    \item \verb|\cref{ch:inleiding}| → \cref{ch:inleiding}
    \item \verb|\Cref{ch:theorie}| → \Cref{ch:theorie} (aan begin van zin)
\end{itemize}
